\section{一个 fd 背后有几个对象}

文件描述符不是文件，也不是内核对象的地址。它只是某个进程文件表中的局部
整数索引；表项持有强引用，强引用指向 \code{FileDescription}，后者才保存
一次打开实例的共享偏移、状态标志与底层 I/O 端点。把这四层分开以后，复制、
关闭、执行时关闭、限额和最终析构才有各自明确的位置。

\topic{历史问题：为什么 Unix 没有把路径直接交给每次读写}
早期文件接口必须同时解决三个彼此不同的问题。第一，目录中的名字可能被
重命名，进程却仍要继续访问已经打开的对象；第二，两次独立打开同一路径应有
独立读写位置；第三，父子进程或复制操作又需要有意共享一次打开的读写位置。
如果每次 \code{read} 都重新解析路径，这三种语义无法同时成立，路径查找成本
也会落入每次 I/O 的热路径。

Unix 系接口因此逐渐形成三层模型：

\begin{enumerate}
  \item 进程中的 descriptor table 把小整数映射为引用；
  \item system-wide open file description 表示一次打开实例，保存偏移和
        file status flags；
  \item inode 或 vnode 表示命名空间与文件系统中的底层对象。
\end{enumerate}

早期的八槽 \code{IoDescriptor} 已经让控制台、管道和文件使用同一组系统
调用，却把“进程索引”“打开实例”和“设备载荷”压在一个固定数组元素里。
这种表示能证明接口形状，不能正确表达以下演进：

\begin{itemize}
  \item \code{dup} 必须增加表项，但不能复制文件偏移；
  \item 两次 \code{open} 必须得到两个打开实例，即使路径相同；
  \item \code{close-on-exec} 属于某一个 fd，不能污染共享打开实例；
  \item 表容量不应由八个编译期槽位决定；
  \item 最后一个引用离开时，文件、管道端点或设备载荷必须恰好析构一次。
\end{itemize}

只把数组从 8 改成 4096 会扩大旧表示，却仍然无法表达共享偏移和最后关闭。
代码先建立通用 \code{KernelObject} 生命周期，再用
\code{FileDescription} 适配已有 I/O，最后让每个 Process 持有动态
\code{FileTable}。

\topic{关闭 fd 时哪一层真正销毁对象}\par
\begin{lstlisting}[caption={Process、FileTable 与共享 FileDescription 的所有权图}]
Process
  |
  +-- FileTable
        |
        +-- fd 0  -- handle --+
        +-- fd 1  -- handle   |
        +-- fd 64 -- handle --+--> KernelObject<FileDescription>
                                      |
                                      +-- shared offset/status flags
                                      +-- console / pipe / legacy-fs payload
\end{lstlisting}

图中是 256 MiB/64 GiB 档的例子，fd 0 与 fd 64 可以指向同一代对象。
数字 64 只在这个 Process 内有效；
另一个 Process 的 fd 64 可以为空，也可以指向完全不同的对象。
\code{KernelObjectHandle} 由所有权容器持有，\code{KernelObjectReference}
则是一次内核操作期间的临时强引用。二者都不向用户态暴露地址。

当前 \code{FileDescription} 直接保存 legacy 文件系统的
\code{FileSystemHandle}，还没有独立 Vnode。这里先确定
打开实例与 fd 的语义，下一章再在它下面接入 Path、Mount 与 Vnode。若同时
同时替换描述符表、命名空间和磁盘格式，失败时将无法判断错误属于哪层。

\section{KernelObject 怎样保存类型和引用}
每个动态对象由一个 64 字节对齐的私有头和紧随其后的类型载荷组成：

\begin{osgridlongtable}{L{0.29\textwidth}L{0.17\textwidth}L{0.42\textwidth}}
字段 & 宽度 & 作用\\ \hline
\code{previous/next} & 各 64 位 & 活动对象侵入链，不另行分配节点\\ \hline
\code{type/state} & 各 64 位 & 类型门禁与 Active/Finalizing 状态机\\ \hline
\code{variant} & 64 位 & 记录 FileDescription 的具体种类\\ \hline
\code{generation} & 64 位 & 区分同址堆内存的不同生命代次\\ \hline
\code{strong_reference_count} & 64 位 & 当前强引用总数\\ \hline
\code{allocation_size_bytes} & 64 位 & 精确归还整个堆块\\ \hline
\code{payload_size_bytes} & 64 位 & 检查载荷边界\\ \hline
\code{finalize_operation/context} & 各 64 位 & 最后释放时的类型析构回调\\ \hline
\code{operation_lock} & 实现定义 & 串行化该对象的可变语义状态\\ \hline
\end{osgridlongtable}

所有位数都由 \code{uint64\_t} 等固定宽度类型表达。代码不把宿主平台的
\code{long}、\code{size\_t} 或指针宽度当成磁盘、ABI 或统计规格。对象地址
本身只存在于 x86-64 内核内部；一旦越过用户 ABI，用户只看到 64 位无符号
描述符参数和 64 位有符号结果。

创建采用“先准备、后发布”：

\begin{enumerate}
  \item 检查管理器、对象类型、载荷大小与溢出；
  \item 从 512 KiB 可回收 kernel heap 取得头加载荷的连续块；
  \item 把载荷逐字节清零，构造头部，但状态仍为 Inactive；
  \item 在管理器锁内取得单调 generation，启动计数 1；
  \item 改为 Active，插入活动链并一次性更新统计；
  \item 最后才把 \code{KernelObjectReference} 交给调用者。
\end{enumerate}

在第 4 步之前发生的任何错误只需释放未发布堆块。对象没有进入活动链，
统计也没有宣称它存在。这个顺序把失败原子性落实成可检查的状态，而不是依靠
调用者猜测“失败到哪一步”。

\topic{地址加 generation：处理内存复用造成的 ABA}
堆释放一个对象后，下一次 best-fit 分配完全可能返回同一虚拟地址。若 handle
只保存地址，旧 handle 会把新对象误认成原对象，这就是一种 ABA：

\begin{lstlisting}[caption={地址复用时由 generation 拒绝旧 handle}]
address P: Object A, generation 41
close last reference -> P is freed
address P: Object B, generation 42
stale handle {P}      -> would incorrectly reach B
safe handle {P, 41}   -> rejected because current generation is 42
\end{lstlisting}

因此 handle 保存 \code{storage\_} 与 \code{generation\_}。取得引用、
读取身份、分离所有权和释放都同时验证：

\[
\begin{aligned}
  &\text{manager matches}\\
  {}\land{}&\text{address is on active list}\\
  {}\land{}&\text{state = Active}\\
  {}\land{}&\text{stored generation = handle generation}.
\end{aligned}
\]

generation 从 1 单调增长，0 永远表示空。到达
\code{UINT64\_MAX} 时创建被明确拒绝，不发生回绕。64 位空间在实验规模下
几乎不会耗尽，但“很难发生”不能替代状态定义；拒绝回绕也使随机模型不必为
代次碰撞留出隐藏例外。

\topic{强引用状态机与最后释放}
\code{ReferenceCounter} 已经处理三条规则：非零启动、零后禁止复活、
上溢拒绝。动态对象把同一操作用于原始 \code{uint64\_t} 状态，
原有类封装继续委托给同一实现，因此动态对象没有复制另一套计数规则。

一次查找的状态变化为：

\[
  R \xrightarrow{\text{AcquireHandle}} R+1
    \xrightarrow{\text{operation}}
  R \xrightarrow{\text{Reference destructor}} R-1.
\]

关闭一个非最后 fd 只把 \(R\) 减一。最后释放则执行更严格的序列：

\begin{enumerate}
  \item 在对象管理器锁内重新验证地址、代次和 Active 状态；
  \item 先检查全局活动引用统计可安全递减，再操作对象计数；
  \item 当结果不是零时，只更新引用统计并返回；
  \item 当结果为零时，将对象改为 Finalizing 并从活动链摘除；
  \item 离开管理器锁，在对象的语义边界执行 finalizer；
  \item 无论类型析构成功与否，都记录结果并归还对象堆块。
\end{enumerate}

finalizer 不能在管理器全局锁内调用。文件系统关闭或管道端点关闭可能进入
其他模块；若全局锁仍被持有，就会把对象图的任意下游锁塞进管理器锁序，
以后极易形成死锁。先从活动集合摘除又保证 finalizer 期间不能取得新引用。

析构失败不会让已经到零的对象复活。它会增加
\code{failed\_finalization\_count}，让调用路径返回明确失败，同时仍释放无主
存储。这样既不制造悬空“僵尸对象”，也不把设备或文件系统错误从证据中抹掉。

\topic{Reference 与 Handle 为什么是两种类型}
\code{KernelObjectReference} 不可复制、可以移动，析构时自动释放。它适合
\code{read}、\code{write} 或状态读取这样的有界调用。若函数提前返回，
RAII 仍会归还临时引用。

\code{KernelObjectHandle} 则只能由 \code{FileTable} 和对象管理器操作。
安装时 \code{DetachReference} 把一份活跃引用无损转换为 handle：计数不变，
只是所有权从局部 RAII 变量转给长期容器。查找时
\code{AcquireHandle} 新增临时引用。关闭时 \code{ReleaseHandle} 消耗表中
那份长期引用。

这个区分排除了两个常见错误：

\begin{itemize}
  \item 安装后局部 reference 析构，意外把表中对象提前释放；
  \item 表项直接保存可自动析构的 C++ 对象，分块清零、移动或销毁时发生
        隐式复制与双重释放。
\end{itemize}

\section{打开实例、文件描述符与进程文件表}

\code{FileDescriptionStorage} 的逻辑内容如下：

\begin{osgridlongtable}{L{0.31\textwidth}L{0.57\textwidth}}
字段 & 语义\\ \hline
\code{kind} & ConsoleInput、ConsoleOutput、ConsoleError、RegularFile、
Directory、PipeReader 或 PipeWriter\\ \hline
\code{file_status_flags} & 可读、可写等共享打开状态\\ \hline
\code{console_input} & 指向 PS/2 字符进入的有界控制台 FIFO\\ \hline
\code{device_write_operation/context} & 注入串口或控制台写操作，不把设备
实现硬编码进对象层\\ \hline
\code{pipe} & 指向现有 64 字节有界 Pipe 及相应端点\\ \hline
\code{file_system} & 指向已挂载的 legacy 文件系统实例\\ \hline
\code{file_system_handle} & inode、类型与当前共享偏移\\ \hline
\end{osgridlongtable}

创建请求必须与 kind 匹配。例如 ConsoleInput 必须提供输入对象且具有可读
状态，ConsoleOutput 必须有写回调且具有可写状态，RegularFile 与 Directory
必须提供文件系统，PipeReader 与 PipeWriter 必须提供相应方向。非法组合在
对象发布前被拒绝。

每个对象自带 \code{operation\_lock}。文件读写、目录枚举和快照都在这个锁
下观察或推进共享偏移。由此可得：

\[
\begin{aligned}
\operatorname{dup}(fd_a) &= fd_b,\\
D(fd_a) &= D(fd_b),\\
\operatorname{read}(fd_a,n) &: offset \leftarrow offset+n,\\
\operatorname{read}(fd_b,m) &: \text{从更新后的 offset 继续}.
\end{aligned}
\]

两次独立 \code{open} 则创建两个 generation 不同的 FileDescription，
两份 \code{FileSystemHandle.offset\_bytes} 都从零开始。路径相同不意味着
打开实例相同。

\topic{共享状态标志与局部描述符标志}
两组名字相似的 flags 必须分开：

\begin{osgridlongtable}{L{0.25\textwidth}L{0.29\textwidth}L{0.34\textwidth}}
类别 & 存储位置 & 复制后的关系\\ \hline
file status flags & FileDescription & 共享；描述一次打开的读写能力\\ \hline
descriptor flags & FileTableEntry & 独立；当前只有 close-on-exec\\ \hline
\end{osgridlongtable}

若把 close-on-exec 放入 FileDescription，那么给复制出来的 fd 设置该位会
连带修改源 fd；下一次 exec 将同时关闭两者。这与 fd 局部语义冲突。
\code{Duplicate} 接受目标 descriptor flags，源表项完全不变。

当前系统还没有完整 \code{execve}，但 \code{CloseOnExec} 已经实现并通过
宿主单元测试。先确定位置和批量关闭语义，接入真正 exec 时便不需要
重塑对象所有权。

\section{FileTable 怎样分配和复制槽位}
每个 Process 持有独立 FileTable。表不预分配 4096 个表项，而是按 64 个
descriptor 一块延迟增长：

\[
  base(fd) = fd - (fd \bmod 64),\qquad index(fd)=fd\bmod 64.
\]

分块按 \code{base\_descriptor} 升序链接。每个表项含一个 16 字节
handle 和一个 64 位 descriptor flags，64 个表项连同块头并考虑 64 字节
对齐约占 1.6 KiB。只使用 fd 0、1、2、3 的进程因此只需要一个分块；复制到
fd 64 才创建第二块。关闭表项不会立即释放中间空块，Process 销毁时统一
释放，使热路径更简单且统计可预测。

三档 hard limit 对应不同运行环境：

\begin{osgridlongtable}{L{0.24\textwidth}L{0.20\textwidth}L{0.44\textwidth}}
配置 & hard limit & 用途\\ \hline
最小调度事务 & 64 & 验证小容量拒绝与失败原子性\\ \hline
functional QEMU & 256 & 运行 Shell、文件、管道与用户态 fd 测试\\ \hline
capacity QEMU & 4096 & 在 64 GiB 主规格上证明表结构不被教学小值锁死\\ \hline
\end{osgridlongtable}

4096 是当前实现的能力上限，不表示已经复制 Linux 的动态上限。
数据结构、计数、参数和系统调用均使用 64 位固定宽度类型，未来提高上限无需
改变 ABI；真正提高前仍要重新评估 heap、每进程内存预算与遍历成本。

\topic{soft limit 与 hard limit 的不同职责}
hard limit 在 FileTable 初始化时确定，且不能超过 4096。soft limit 可由
当前进程调整，但必须满足：

\[
  1 \le soft \le hard \le 4096.
\]

soft limit 限制新安装可取得的 fd 数值范围，不会因为降低限额而强制关闭已经
存在的高编号 fd。这样限额调整不会产生隐式 I/O 副作用。查找和关闭已有 fd
按 hard limit 验证，安装与复制按当前 soft limit 搜索。

若最小目标 descriptor 已经不小于 soft limit，操作返回
\code{SoftLimitExceeded}。用户 ABI 将其映射为 \code{-24}。底层对象引用、
原表项、分块集合与计数都保持原值。

\topic{为什么分配必须返回最低可用 fd}
用户程序长期依赖“从给定下界开始的最低可用描述符”语义。它不只是方便打印：
标准输入、输出和错误固定为 0、1、2；关闭 fd 4 后再次打开，若没有更小空位，
就应稳定复用 4。确定性也使整机证据能够检查确切数字。

\code{FindLowestAvailableDescriptor} 从调用方下界开始逐块扫描。遇到未分配
块时，它同时返回该块的 base，调用方先建立块，再重新搜索。它不能把“缺块”
直接视为 64 个永久空槽后立刻发布，因为分配新块需要 heap，而表锁内不应嵌套
执行复杂分配。

\topic{安装的两阶段提交}
\code{Install} 的关键循环可以写成以下事务：

\begin{lstlisting}[caption={FileTable 安装的两阶段提交}]
lock table
  validate soft limit
  find lowest slot
  if its chunk is missing:
      unlock table
      allocate and zero candidate chunk
      lock table, publish chunk if still absent, unlock
      retry search
  verify slot still empty
  detach Reference -> Handle
  publish handle + fd flags
  update counts
unlock table
\end{lstlisting}

\code{EnsureChunk} 也采用候选块。若另一个执行流已发布相同 base，候选块会被
释放，现有块继续使用。当前 BSP 还没有多核并发，但把竞争语义写进接口能防止
以后通过加 AP 才发现安装路径天生不可线性化。

最重要的提交点是 \code{DetachReference}。在它之前，失败时调用者仍拥有
reference；在它之后，handle 与统计在同一表锁临界区内发布。用户永远看不到
只有 flags 没有对象，或只有对象却未计数的半表项。

\topic{查找、复制、关闭和销毁的完整路径}
四个动作围绕同一条所有权主线展开：先在锁内改变表项和引用，再把可能触发
复杂析构的工作移到锁外，从而同时守住可见性与生命周期。

\topic{Lookup}

表锁保护 chunk 与 entry 的稳定性。在锁内验证 fd、找到非空 handle，再调用
\code{AcquireHandle}。取得临时强引用后才释放表锁，因此并发 close 只能
释放表的那份引用，正在进行的 I/O 仍拥有自己的引用。

\topic{Duplicate}

先 Lookup 源 fd，得到同一 generation 的临时引用；再按目标下界 Install。
成功时临时引用被转移为新表项，因此对象计数净增 1。安装失败时 RAII 析构
临时引用，计数恢复，源 fd 不变。

\topic{Close}

先在表锁内清空 entry 并递减 active descriptor，再在锁外释放 handle。
最后引用可能触发文件系统或管道 finalizer，所以不能在表锁内执行。表项一旦
清空，新 Lookup 就会失败；已经取得的临时引用仍可安全完成。

\topic{Destroy}

Process 回收时逐项释放所有非空 handle，再释放全部 chunk，最后检查活动
descriptor 与 chunk 归零。若任一 finalizer 失败，销毁返回失败并记录证据，
但不会故意保留已失去所有者的表内存。

\topic{锁顺序与目前的并发边界}
当前主要锁层次为：

\[
  \text{FileTable lock}
  \rightarrow \text{KernelObjectManager lock}
  \rightarrow \text{temporary reference acquired}.
\]

实际 I/O 在离开 FileTable 锁后，通过对象的 operation lock 进入控制台、
Pipe 或 FileSystem。最后析构也在对象管理器和表锁之外进入下游模块。统计有
独立短临界区，不在每个读写字节打印日志。

当前内核仍是单 BSP，但 PIT IRQ 可在系统调用期间到达，调度也可在返回前
切换 Thread。取得 FileTable 锁以后先完成表项替换，释放它以后才可能进入
对象 finalizer，这个顺序避免两把锁互相等待。当前没有跨核 RCU、无锁 fd
查找或远端 close；这些能力还需要 AP、IPI 与内存回收协议。

\topic{既有 I/O 如何无旁路迁移}
迁移后的通用路径统一为：

\[
\begin{aligned}
\text{fd}
&\rightarrow \text{FileTable::Lookup}
\rightarrow \text{KernelObjectReference}\\
&\rightarrow \text{FileDescriptionManager}
\rightarrow \text{backend}.
\end{aligned}
\]

\begin{osgridlongtable}{L{0.29\textwidth}L{0.29\textwidth}L{0.30\textwidth}}
端点 & 读路径 & 写路径\\ \hline
ConsoleInput & 读取 256 字节 FIFO，可返回 WouldBlock & 拒绝\\ \hline
ConsoleOutput/Error & 拒绝 & 调用注入的设备写函数\\ \hline
RegularFile & legacy-fs Read 并推进共享 offset & Write 并推进共享 offset\\ \hline
Directory & ReadDirectory 并推进目录位置 & 拒绝\\ \hline
PipeReader & Pipe TryRead、EOF 与唤醒语义 & 拒绝\\ \hline
PipeWriter & 拒绝 & Pipe TryWrite、背压与 broken pipe\\ \hline
\end{osgridlongtable}

原先专用的文件、管道和控制台系统调用仍保留 ABI 兼容，但它们最终汇入同一
对象路径。没有旧 \code{IoDescriptor} 旁路，也没有为 PID 4 的证明写特殊
内核读取函数。这样历史 QEMU 路径与新能力同时检验迁移是否完整。

\section{从系统调用到整机测试}
原生 \code{SYSCALL} 与 \code{INT 0x80} 已经共用
\code{UserContext}，fd 操作只在分发层增加：

\begin{osgridlongtable}{L{0.10\textwidth}L{0.31\textwidth}L{0.35\textwidth}L{0.12\textwidth}}
编号 & 名称 & 参数 & 成功结果\\ \hline
23 & DuplicateDescriptor & RDI=源 fd，RSI=最低目标，RDX=目标 flags & 新 fd\\ \hline
24 & GetDescriptorFlags & RDI=fd & flags\\ \hline
25 & SetDescriptorFlags & RDI=fd，RSI=flags & 0\\ \hline
26 & SetDescriptorSoftLimit & RDI=新 soft limit & 0\\ \hline
27 & GetDescriptorSoftLimit & 无 & soft limit\\ \hline
28 & GetDescriptorHardLimit & 无 & hard limit\\ \hline
\end{osgridlongtable}

RAX 在进入前保存系统调用号，返回时保存 64 位结果。负数沿既有 ABI 表示错误：
\code{-12} 是非法 fd，\code{-24} 是限额拒绝，\code{-25} 是内核对象失败。
用户包装不读取平台相关的 \code{errno}，也不依赖宿主 C 库。

\topic{对象模型变化为什么没有改写汇编入口}
本阶段仍然真实经过 Intel 语法汇编，而不是绕过汇编。以
\code{DuplicateDescriptor} 为例，控制流是：

\begin{enumerate}
  \item freestanding C++ 包装把 23 放入 RAX，把源 fd、下界和 flags 分别放入
        RDI、RSI、RDX；
  \item \filepath{source/user/src/system\_call.asm} 执行 \code{syscall}；
  \item CPU 按 STAR/LSTAR/FMASK 与 EFER 配置进入 Ring 0，把用户 RIP 放入
        RCX、RFLAGS 放入 R11；
  \item \filepath{arch/architecture.asm} 执行 \code{swapgs}，从 CpuLocal
        取得可信栈，构造统一 176 字节 UserContext；
  \item C++ dispatcher 读取 RAX/RDI/RSI/RDX，调用当前 Process 的 FileTable；
  \item 结果写回现场 RAX，经合法性检查选择 \code{sysretq} 或
        \code{iretq} 返回。
\end{enumerate}

没有增加新的寄存器，也没有改变汇编帧偏移。硬件特权转换与对象操作已经在
UserContext 边界分开，因此这里只需修改 C++ dispatcher 和文件表。
若新增一个 fd 操作就必须修改 SWAPGS 或返回帧，说明 ABI 层次没有真正稳定。

\topic{PID 4 的整机语义证明}
调度 worker 在真实 Ring 3 中执行以下序列：

\begin{enumerate}
  \item 创建并写入 \filepath{/fdv14.bin}，关闭写 fd；
  \item 以读模式打开为 fd 3，查询 hard limit 后选择复制下界 \(m\)，并设置
        close-on-exec；
  \item 再次独立打开同一路径，最低空位应为 fd 4；
  \item 从 fd 3 读取前三字节，再从复制 fd 读取后续三字节；
  \item 从独立 fd 4 读取前三字节，证明它没有共享前两次读取推进的偏移；
  \item 比较源与复制项的 descriptor flags，修改源 flags 后再次读取；
  \item 把 soft limit 降为 \(m\)，要求从最低 \(m\) 复制明确得到 \code{-24}；
  \item 恢复到查询所得 hard limit，关闭 fd 4，再次打开并要求最低空位复用 4；
  \item 关闭全部描述符，打印一次有界成功标记。
\end{enumerate}

\(m\) 在 256 MiB 与 64 GiB 档为 64，在 hard limit 本身只有 64 的兼容档
为 8。后者不能使用 minimum 64，因为合法候选区间
\([64,64)\) 为空。选择 8 不是跳过语义：它仍在连续低编号之外制造空洞，
并验证最低编号选择、共享 offset、独立 flags、限额拒绝与失败原子性。

载荷是 8 字节；整个 PID 4 文件描述模型累计读取 9 字节、写入 8 字节。
历史 Shell、管道和文件系统仍执行自己的路径，因此最终管理器统计会包含更多
对象与字节，检查脚本不会把某次局部事务误当成全系统总量。

\topic{统计为何在阶段结束时集中打印}
fd 查找、字符读写和引用增减都是热路径。逐次串口输出会改变 PIT 调度时序，
淹没真正异常，还可能让 QEMU 超时。内核只在这组操作结束后打印汇总：

\begin{lstlisting}[caption={functional 运行的一次对象统计}]
[OS][KERNEL] OBJECT_ACTIVE_COUNT=0x0000000000000000
[OS][KERNEL] OBJECT_ACTIVE_REFERENCES=0x0000000000000000
[OS][KERNEL] OBJECT_CREATIONS=0x0000000000000017
[OS][KERNEL] OBJECT_DESTRUCTIONS=0x0000000000000017
[OS][KERNEL] FILE_DESCRIPTION_FINALIZATIONS=0x0000000000000017
[OS][KERNEL] FILE_DESCRIPTION_FAILED_FINALIZATIONS=0x0000000000000000
[OS][KERNEL] FILE_TABLE_HARD_LIMIT=0x0000000000000100
[OS][KERNEL] FILE_TABLE_CHUNK_ALLOCATIONS=0x0000000000000005
[OS][KERNEL] FILE_TABLE_CHUNK_RELEASES=0x0000000000000005
[OS][KERNEL] PROCESS_RESOURCE_VALIDATION=0x0000000000000001
[OS][KERNEL] RUNTIME_STATE_VALIDATION=0x0000000000000001
[OS][KERNEL] SMOKE_STATE_VALIDATION=0x0000000000000001
\end{lstlisting}

同一 QEMU 日志还打印 \code{TIMER\_TICKS}、
\code{MONOTONIC\_MILLISECONDS} 与调度 tick，给来宾内部时间提供可比较
坐标；宿主运行器另记录整体阶段耗时和有界超时。创建等于析构、分块申请等于
释放、活动对象与强引用为零、finalization failure 为零，四者联合证明生命
周期；只检查 \code{READY} 无法排除泄漏。

\topic{四层测试怎样互补}\par
\begin{osgridlongtable}{L{0.25\textwidth}L{0.28\textwidth}L{0.35\textwidth}}
层次 & 代表测试 & 主要排除的问题\\ \hline
单元 & \filepath{file\_table\_test.cpp} & 引用转移、共享 generation、独立
fd flags、最低复用、close-on-exec 与销毁\\ \hline
集成 & \filepath{file\_description\_lifecycle\_test.cpp} & legacy-fs、
管道和对象 finalizer 的跨模块生命周期\\ \hline
容量 & \filepath{file\_table\_capacity\_test.cpp} & 4096 fd、64 个分块、
耗尽失败原子性与完全释放\\ \hline
随机 & \filepath{file\_table\_randomized\_test.cpp} & 固定种子 100000 步
模型对照，每一步核对占用、generation、flags、引用与统计\\ \hline
整机 & functional/capacity QEMU & Ring 3 ABI、真实汇编入口、调度、
文件系统、heap 与最终全局守恒\\ \hline
\end{osgridlongtable}

随机模型不把实现本身当预言机。宿主侧以独立数组记录每个 fd 是否占用、
对应 generation 与 flags，再把 Install、Duplicate、Lookup、Close 和限额
操作同时作用于模型和真实实现。每一步都比较整个可见状态，结束时要求所有
对象与堆分配归零。

对象和 FileTable 测试加入后，早期回归与实体学习图连通性门禁仍然保留。
历史固定槽实现的专用测试随生产类型一同删除；继续测试已经不存在的数据结构
无法保护当前实现。

\topic{增量构建为什么必须依赖生成文件}
用户 worker 先链接成 User ELF，再由
\filepath{user\_images.o} 嵌入 Kernel，随后依次生成
\filepath{kernel.elf}、\filepath{kernel.payload.elf} 和磁盘镜像：

\begin{lstlisting}[caption={用户程序进入启动磁盘的真实依赖链}]
scheduler_worker.cpp
  -> user_scheduler_worker.elf
  -> user_images.o
  -> kernel.elf
  -> kernel.payload.elf
  -> boot_disk.img
\end{lstlisting}

CMake 的自定义 target 是 phony 节点。它适合声明“生产者必须先运行”的顺序，
却不能单独把生成文件的时间戳传播给下游自定义命令。若 Kernel 链接只依赖
\code{os\_kernel\_user\_images\_object} target，User ELF 改变后
\filepath{user\_images.o} 可以更新，而旧 \filepath{kernel.elf} 仍被认为
最新，最终表现为源码已经修改但 QEMU 仍运行旧 worker。

因此 Kernel 链接同时依赖 target 和
\filepath{architecture.o}/\filepath{user\_images.o} 两个真实输出。前者保证
跨目录构建顺序，后者保证时间戳使链接失效。磁盘边界审计则只检查实际写盘的
\filepath{kernel.payload.elf}；完整 \filepath{kernel.elf} 保留 DWARF 供
GDB 使用，其非加载调试段不属于 Stage 1 暂存载荷。这个区分同时修复了
“旧镜像假失败”和“把调试文件长度当启动载荷长度”两类工程错误。

\topic{沿目录走读对象表与文件表}\par
\begin{osgridlongtable}{L{0.43\textwidth}L{0.45\textwidth}}
文件 & 阅读时要回答的问题\\ \hline
\filepath{foundation/reference\_counter.hpp/.cpp} &
原始计数操作怎样与旧类封装共享零后不可复活语义？\\ \hline
\filepath{kernel/object/kernel\_object.hpp/.cpp} &
创建、handle/reference 转换、generation 检查与最后析构的提交点在哪里？\\ \hline
\filepath{kernel/io/file\_description.hpp/.cpp} &
哪些状态属于共享打开实例，控制台、文件、目录和管道如何统一？\\ \hline
\filepath{kernel/io/file\_table.hpp/.cpp} &
64 槽分块、最低 fd、软硬限额和两阶段安装怎样保持失败原子性？\\ \hline
\filepath{kernel/process/process\_runtime.hpp/.cpp} &
每个 Process 怎样创建标准 fd、接入 Pipe，并在退出时销毁 FileTable？\\ \hline
\filepath{kernel/user/system\_calls.cpp} &
23--28 号调用怎样验证参数、映射 \code{-24/-25} 并复用双入口？\\ \hline
\filepath{user/src/system\_call.cpp} &
freestanding 包装怎样保持 RAX/RDI/RSI/RDX ABI？\\ \hline
\filepath{user/programs/scheduler\_worker.cpp} &
共享偏移、独立打开、fd flags、限额和最低复用怎样由 Ring 3 证明？\\ \hline
\filepath{tests/randomized/kernel/file\_table\_randomized\_test.cpp} &
独立模型在十万步中保存哪些状态，何时检查引用守恒？\\ \hline
\filepath{tools/os\_tools/qemu\_runner.py} &
哪些串口 marker 被精确解析，超时怎样终止而不永久卡住？\\ \hline
\end{osgridlongtable}

推荐先从 unit test 推导公开语义，再读 FileTable；随后向下进入
KernelObject，向上进入 ProcessRuntime 与系统调用。最后阅读 PID 4 与 QEMU
解析器，检查模型承诺是否真的穿过硬件入口到达整机证据。

\topic{当前实现范围}
当前还没有 Vnode、Path、Mount、memfs 或新磁盘格式；legacy-fs handle 仍直接
存在于 FileDescription。没有 fork，因此还没有跨 Process 继承表；没有完整
exec，因此 close-on-exec 只在对象层测试。没有 seek、append、非阻塞 flag
或 advisory lock。对象类型当前只有 FileDescription，通用 manager 的目的
是让最后一次释放有确定路径，不是提前制造庞大类层次。

FileTable 的空 chunk 只在 Destroy 时释放；4096 是当前 hard limit，
不是 Linux 兼容声明。KernelObject finalizer 仍是同步调用，也没有弱引用、
循环引用检测、RCU 延迟释放或多核缓存。把这些边界写出，下一章才能
聚焦 VFS，而不是把尚未实现的能力藏在“类似 Linux”的模糊表述中。

运行 fd 测试后，进程局部名字与共享打开实例应当仍然分离：复制 fd 会共享
offset，两次独立 open 不共享 offset，close-on-exec 只改变一个表项。地址
复用由 generation 检出，最后一个强引用触发一次 finalizer。容量与随机测试
还要在每次失败后检查原表项、引用数和已分配 chunk 没有变化，并在销毁后回到
零。
